\section{行政协议（行政合同）}
\begin{enumerate}
  \item 官民双方签的就是行政合同。\bs{为了对外实施行政管理，具有行政法上权利义务内容}。
  \item 不能官告民、官不能提反诉！！民只能提行诉。
  \item 种类。
        \begin{enumerate}
          \item 特许经营（公共事业承包给私人）。
          \item \zd{土地、房屋的征收补偿}。
          \item 自然资源的出让。
          \item 社会资本合作/PPP。
        \end{enumerate}
  \item 审合法性和合约性。可以调解。
  \item 非诉执行申请。签合同 + 管合同。
  \item 管辖。
        \begin{enumerate}
          \item 推定管辖。提过民诉 \yjt 不立案 \yjt 又提行诉 \yjt 必须立案。处理实践踢皮球。
          \item 级别、专属管辖不能协商。\al{当事人约定由基层法院管辖 \cuo}。
          \item 约定仲裁条款\zd{无效}。但法律、行政法规或者我国缔结、参加的国际条约\zd{另有规定的除外}。
        \end{enumerate}
  \item 行政协议无效的原因在一审法庭辩论终结前消除的，可以确认行政协议有效。
  \item 法律适用。
        \begin{enumerate}
          \item 诉讼问题民诉（\al{怎么判}），实体问题民法（\al{合同是否无效、可撤销}）。\zd{2015 年往后，新案新办法}。
          \item 例外。一头（\zx{起诉期}）一尾（\zx{判决}）一中间（\zx{举证责任}），一分为二。
                政府\zd{单方变更、解除}！！ \yjt 起诉期、判决、举证 \yjt 走行政，其他走民诉。
                \al{达成 100W 补偿协议，\xhx{区政府}发现面积有误，将补偿\xhx{变更}为 80W}。
          \item \bs{本案的[起诉期限、判决、举证]应当适用什么规定}。
                \begin{enumerate}
                  \item \bs{不依法履行、未按照约定履行行政协议}，xx \zd{参照民事}法律规范确定。
                  \item 对行政机关\zd{变更、解除}行政协议等行政行为起讼的，xx \zd{依照行政}诉讼法及其司法解释确定。
                \end{enumerate}
        \end{enumerate}
  \item 判决。
        \begin{enumerate}
          \item 走行政。
                \begin{enumerate}
                  \item 合法。判决驳回原告诉讼请求；给原告造成损失的，判决被告\zd{予以补偿}。
                  \item 违法。\zx{撤销 \hz 部分撤销，并可以责令重新作出}。
                \end{enumerate}
          \item 走民诉。
                \begin{enumerate}
                  \item 合法。驳回原告诉讼请求；原告遭受损失的，判决被告\zd{予以补偿}。
                  \item 违法。
                        \begin{enumerate}
                          \item 判决被告\zd{继续履行}，并明确继续履行的具体内容。
                          \item 无法履行 \hz 履行无意义。判决\zd{采取相应的补救措施}。
                                原告遭受损失的，判决被告\zd{予以补偿}（要求违约金、定金，支持）。
                          \item 预期违约。原告请求被告\zd{承担违约责任}的，支持。
                        \end{enumerate}
                \end{enumerate}
          \item 也存在转化问题。
                \begin{enumerate}
                  \item 原告以被告\zd{违约}为由请求法院判令其承担违约责任，法院经审理认为行政协议\zd{无效的}，应当向原告\zd{释明}，并根据原告变更后的诉讼请求\zd{判决确认行政协议无效}。
                  \item 因被告的行为造成\zx{行政协议无效的}，法院可以依法判决被告承担赔偿责任。原告经释明后\zx{拒绝变更诉讼请求的}，法院可以\zx{判决驳回其诉讼请求}。
                \end{enumerate}
        \end{enumerate}
\end{enumerate}
